\section{Methodology}

\subsection{Problem Setting}

Consider a code generation model $\pi_\theta$ trained with PPO on execution feedback. Given a problem description $x$, the model generates a code solution $y = (y_1, \ldots, y_T)$ which is executed against test cases. The reward $r(y)$ is typically binary: 1 if all tests pass, 0 otherwise.

Standard PPO applies the reward uniformly across all tokens. This ignores a key observation: not all tokens influence the test outcome.

\subsection{Fine-Grained Optimization}

FGO addresses this by masking non-executed tokens from gradient updates. Given an execution trace $\tau$ collected during test evaluation, we construct a binary mask $m \in \{0, 1\}^T$ where $m_t = 1$ if token $y_t$ corresponds to executed code.

The FGO loss modifies PPO to only update executed tokens, concentrating gradients on tokens that actually influenced the outcome.

\subsection{Mechanism Decomposition}

We decompose FGO into three components:

\textbf{Component 1: Trace Collection (H-M1)}: Python's \texttt{sys.settrace} captures executed line numbers. Falsification: capture rate $< 95\%$.

\textbf{Component 2: Token Classification}: Tokenizer offset\_mapping maps tokens to lines. Falsification: F1 $< 70\%$.

\textbf{Component 3: Gradient Exclusion (H-M2)}: Masked tokens receive zero gradient. Falsification: any non-zero gradient for masked tokens.

\textbf{Component 4: Credit Assignment (H-M3)}: Signal concentration improves learning. Falsification: FGO pass@1 $\leq$ Standard pass@1.

\subsection{Verification Protocol}

\begin{table}[h]
\centering
\caption{Verification protocol with gate criteria.}
\label{tab:verification}
\begin{tabular}{@{}llll@{}}
\toprule
Hypothesis & Type & Gate & Pass Criterion \\
\midrule
H-E1 & Existence & MUST\_WORK & FGO improves across all content types \\
H-M1 & Mechanism & MUST\_WORK & Trace $\geq$95\%, F1 $\geq$70\% \\
H-M2 & Mechanism & MUST\_WORK & Non-executed gradient = 0 \\
H-M3 & Efficiency & SHOULD\_WORK & FGO pass@1 $>$ Standard pass@1 \\
\bottomrule
\end{tabular}
\end{table}
